\subsection{Simu5G and Sionna Prior Integration Limitations}
Simu5G\cite{simu5g} is an open-source 5G system-level simulator built on the OMNeT++ framework. It is developed by a consortium of universities and research centers in Spain, led by the Universitat Politècnica de Catalunya (UPC) and has been used for various research works on 5G networks.

Simu5G use Network Description language (NED) for topology and C++ for behavior of the network and applications. Topology is defined in NED file, parameters are passed through a separate omnetpp.ini file and instantiated via a runtime that reflects on those declarations.
While elegant for purely C++ network research, the framework is very rigid and hard to integrate with other frameworks such as SionnaRT which is based on Python.

Another key limitation of Simu5G is its lack of native support for multi-input multi-output (MIMO) antenna systems.
Recent versions of Simu5G doesn’t support MIMO antenna systems natively and require workarounds to enable it which
add another layer of complexity and brittleness to the framework.

In addition, it use statistical models for channel modeling which are not as accurate as ray-traced models that reflect the physical environment. Propagation loss and delay are modeled stochastically, which can be acceptable
for large-scale system-level analysis but is inadequate for high-accuracy simulations of dense urban or indoor environments.
Despite our previous efforts to integrate SionnaRT with Simu5G for MIMO antenna systems with ray-tracing capabilities, we encountered several limitations that motivated our transition to ns-3 for our 6G simulations.
Transitioning to ns-3 significantly provide multiple advantages including easier integration with other frameworks such as SionnaRT which is based on Python, native support for MIMO antenna systems and ability to use thirty party libraries like 5G LENA NR for 5G NR simulations.